\section{仿射联络}\label{sec:affine-connections}

若 $X$ 和 $Y$ 是 Euclidean 空间中的向量场, 可以定义 $Y$ 沿 $X$ 的\textbf{方向导数} $\nabla_XY$. 然而, 这一定义用到了 Cartesian 坐标的存在性, 而一般流形上不再有这种坐标. 为克服这个困难, 必须引入更多结构.

\begin{RiemannDefinition}\label{def:affine-connection}
设 $M$ 是可微流形. $M$ 上的\textbf{仿射联络}是一个映射 $\nabla:\mathfrak{X}(M)\times\mathfrak{X}(M)\to\mathfrak{X}(M)$, 对所有 $X,Y,Z\in\mathfrak{X}(M)$ 和 $f,g\in C^\infty(M,\mathbb{R})$ 满足
\begin{enumerate}[label=(\roman*),leftmargin=*]
\item $\nabla_{fX+gY}Z=f\nabla_XZ+g\nabla_YZ$;
\item $\nabla_X(Y+Z)=\nabla_XY+\nabla_XZ$;
\item $\nabla_X(fY)=(X\cdot f)Y+f\nabla_XY$
\end{enumerate}
(记 $\nabla_XY:=\nabla(X,Y)$).
\end{RiemannDefinition}

向量场 $\nabla_XY$ 有时称为 $Y$ 沿 $X$ 的\textbf{协变导数}.

\begin{RiemannProposition}\label{prop:local-expression-of-affine-connection}
设 $\nabla$ 是 $M$ 上的仿射联络, $X,Y\in\mathfrak{X}(M)$ 且 $p\in M$. 则 $(\nabla_XY)_p\in T_pM$ 只依赖于 $X_p$ 和 $Y$ 沿一条在 $p$ 处与 $X$ 相切的曲线的取值. 此外, 若 $x:W\to\mathbb{R}^n$ 是某个开集 $W\subset M$ 上的局部坐标, 且在该开集上
\[
  X=\sum_{i=1}^n X^i\frac{\partial}{\partial x^i},\qquad
  Y=\sum_{i=1}^n Y^i\frac{\partial}{\partial x^i},
\]
则
\begin{equation}\label{eq:local-expression-of-affine-connection}
  \nabla_XY=\sum_{i=1}^n
  \left(X\cdot Y^i+\sum_{j,k=1}^n\Gamma_{jk}^i X^jY^k\right)
  \frac{\partial}{\partial x^i}
\end{equation}
其中 $n^3$ 个可微函数 $\Gamma_{jk}^i:W\to\mathbb{R}$ 称为 \textbf{Christoffel 符号}, 由下式定义:
\begin{equation}\label{eq:definition-of-christoffel-symbols}
  \nabla_{\frac{\partial}{\partial x^j}}\frac{\partial}{\partial x^k}
  =\sum_{i=1}^n\Gamma_{jk}^i\frac{\partial}{\partial x^i}.
\end{equation}
\end{RiemannProposition}

\begin{RiemannProof}
容易证明仿射联络是\textbf{局部的}, 即若 $X,Y\in\mathfrak{X}(M)$ 在某个开集 $W\subset M$ 上与 $\widetilde X,\widetilde Y\in\mathfrak{X}(M)$ 一致, 则 $\nabla_XY=\nabla_{\widetilde X}\widetilde Y$ 在 $W$ 上成立 [见 Exercise~\ref{ex:affine-connections}(\ref{ex:locality-and-coordinate-expressions})]. 所以可以计算只定义在 $W$ 上的向量场 $X,Y$ 的 $\nabla_XY$. 设 $W$ 是局部坐标 $x:W\to\mathbb{R}^n$ 的坐标邻域, 由 \eqref{eq:definition-of-christoffel-symbols} 定义与这些局部坐标对应的 Christoffel 符号. 展开
\[
  \nabla_XY=\nabla_{\left(\sum_{i=1}^nX^i\frac{\partial}{\partial x^i}\right)}
  \left(\sum_{j=1}^nY^j\frac{\partial}{\partial x^j}\right)
\]
并应用 Definition~\ref{def:affine-connection} 中列出的性质, 得到 \eqref{eq:local-expression-of-affine-connection}. 此式表明, $(\nabla_XY)_p$ 只依赖于 $X^i(p)$, $Y^i(p)$ 和 $(X\cdot Y^i)(p)$. 此外, $X^i(p)$ 和 $Y^i(p)$ 只依赖于 $X_p$ 和 $Y_p$, 而 $(X\cdot Y^i)(p)=\left.\frac{\mathrm{d}}{\mathrm{d}t}Y^i(c(t))\right|_{t=0}$ 只依赖于 $Y^i$ (或 $Y$) 沿曲线 $c$ 的取值, 该曲线在 $p=c(0)$ 处的切向量为 $X_p$.
\end{RiemannProof}

\begin{RiemannRemark}\label{rem:connection-determined-by-christoffel-symbols}
局部而言, 指定一个坐标邻域上的 Christoffel 符号就唯一确定了仿射联络. 但是, 不同坐标图上的 Christoffel 符号不能独立选取, 因为协变导数在重叠处必须一致.
\end{RiemannRemark}

定义在可微曲线 $c:I\to M$ \textbf{沿线的向量场}是一个可微映射 $V:I\to TM$, 对所有 $t\in I$ 满足 $V(t)\in T_{c(t)}M$. 一个显然的例子是切向量 $\dot c(t)$. 若 $V$ 是定义在 $\dot c\ne0$ 的可微曲线 $c:I\to M$ 沿线的向量场, 则其沿 $c$ 的\textbf{协变导数}是定义在 $c$ 沿线并由下式给出的向量场:
\[
  \frac{DV}{\mathrm{d}t}(t):=\nabla_{\dot c(t)}V=(\nabla_XY)_{c(t)}
\]
其中 $X,Y\in\mathfrak{X}(M)$ 为任意满足 $X_{c(t)}=\dot c(t)$ 的向量场, 且对某个 $\varepsilon>0$ 和 $s\in(t-\varepsilon,t+\varepsilon)$ 有 $Y_{c(s)}=V(s)$. 注意, 只要 $\dot c(t)\ne0$, 这样的延拓就总存在. Proposition~\ref{prop:local-expression-of-affine-connection} 保证 $(\nabla_XY)_{c(t)}$ 不依赖于 $X,Y$ 的选取. 事实上, 若在局部坐标 $x:W\to\mathbb{R}^n$ 下有 $x^i(t):=x^i(c(t))$, 且
\[
  V(t)=\sum_{i=1}^n V^i(t)\left(\frac{\partial}{\partial x^i}\right)_{c(t)},
\]
则
\[
  \frac{DV}{\mathrm{d}t}(t)=\sum_{i=1}^n
  \left(\dot V^i(t)+\sum_{j,k=1}^n\Gamma_{jk}^i(c(t))\dot x^j(t)V^k(t)\right)
  \left(\frac{\partial}{\partial x^i}\right)_{c(t)}.
\]

\begin{RiemannDefinition}\label{def:parallel-vector-field-and-geodesic}
定义在曲线 $c:I\to M$ 沿线的向量场 $V$ 称为\textbf{沿 $c$ 平行}, 若对所有 $t\in I$, 有
\[
  \frac{DV}{\mathrm{d}t}(t)=0.
\]
若 $\dot c$ 沿 $c$ 平行, 即对所有 $t\in I$ 有
\[
  \frac{D\dot c}{\mathrm{d}t}(t)=0,
\]
则称曲线 $c$ 为联络 $\nabla$ 的\textbf{测地线}.
\end{RiemannDefinition}

在局部坐标 $x:W\to\mathbb{R}^n$ 下, $V$ 沿 $c$ 平行的条件写为
\begin{equation}\label{eq:parallel-transport-in-local-coordinates}
  \dot V^i+\sum_{j,k=1}^n\Gamma_{jk}^i\dot x^jV^k=0\quad(i=1,\ldots,n).
\end{equation}
这是关于 $V$ 的各分量的一阶线性常微分方程组. 由 Picard–Lindelöf 定理以及线性常微分方程的全局存在定理 \RGcite{ch03-Arn92}, 给定曲线 $c:I\to M$, 点 $p\in c(I)$ 和向量 $v\in T_pM$, 存在唯一的沿 $c$ 平行的向量场 $V:I\to TM$, 使得 $V(0)=v$, 称为 $v$ 沿 $c$ 的\textbf{平行移动}.

此外, 测地线方程为
\begin{equation}\label{eq:geodesic-equations-in-local-coordinates}
  \ddot x^i+\sum_{j,k=1}^n\Gamma_{jk}^i\dot x^j\dot x^k=0\quad(i=1,\ldots,n).
\end{equation}
这是关于 $c(t)$ 的各坐标的二阶 (非线性) 常微分方程组. 因此, 由 Picard–Lindelöf 定理, 给定点 $p\in M$ 和向量 $v\in T_pM$, 存在唯一的测地线 $c:I\to M$, 定义在包含 $0\in I$ 的极大开区间 $I$ 上, 满足 $c(0)=p$ 和 $\dot c(0)=v$.

现在定义仿射联络 $\nabla$ 的挠率. 为此, 注意在局部坐标 $x:W\to\mathbb{R}^n$ 下, 有
\begin{align*}
  \nabla_XY-\nabla_YX
  &=\sum_{i=1}^n\left(X\cdot Y^i-Y\cdot X^i
    +\sum_{j,k=1}^n\Gamma_{jk}^i(X^jY^k-Y^jX^k)\right)\frac{\partial}{\partial x^i}\\
  &=[X,Y]+\sum_{i,j,k=1}^n(\Gamma_{jk}^i-\Gamma_{kj}^i)X^jY^k\frac{\partial}{\partial x^i}.
\end{align*}

\begin{RiemannDefinition}\label{def:torsion-of-affine-connection}
$M$ 上联络 $\nabla$ 的\textbf{挠率算子}是对所有 $X,Y\in\mathfrak{X}(M)$ 由下式给出的算子 $T:\mathfrak{X}(M)\times\mathfrak{X}(M)\to\mathfrak{X}(M)$:
\[
  T(X,Y)=\nabla_XY-\nabla_YX-[X,Y].
\]
若 $T=0$, 则称该联络\textbf{对称}.
\end{RiemannDefinition}

$T(X,Y)$ 的局部表达式表明, $T(X,Y)_p$ 线性依赖于 $X_p$ 和 $Y_p$. 换言之, $T$ 是 $M$ 上一个 $(2,1)$-张量场, 由局部坐标中的下式给出:
\[
  T=\sum_{i,j,k=1}^n(\Gamma_{jk}^i-\Gamma_{kj}^i)
    \,\mathrm{d}x^j\otimes\mathrm{d}x^k\otimes\frac{\partial}{\partial x^i}
\]
(回顾第~\ref{chap:differential-forms}~章的 Remark~\ref{rem:mixed-tensors-and-linear-maps}: 任意 $(2,1)$-张量 $T\in\mathcal{T}^{2,1}(V^*,V)$ 都可自然地与一个双线性映射 $\Phi_T:V^*\times V^*\to V\cong V^{**}$ 等同, 具体地对所有 $v,w\in V$ 和 $\alpha\in V^*$, 有 $\Phi_T(v,w)(\alpha):=T(v,w,\alpha)$).

注意, 联络对称当且仅当对所有 $X,Y\in\mathfrak{X}(M)$ 有 $\nabla_XY-\nabla_YX=[X,Y]$. 在局部坐标中, 联络对称的条件为
\[
  \Gamma_{jk}^i=\Gamma_{kj}^i\quad(i,j,k=1,\ldots,n)
\]
(名称由此而来).

\begin{RiemannExercise}\label{ex:affine-connections}
\begin{enumerate}[label=(\arabic*),ref=\arabic*,leftmargin=*]
\item\label{ex:locality-and-coordinate-expressions}
\begin{enumerate}[label=(\alph*),leftmargin=*]
\item 证明若 $X,Y\in\mathfrak{X}(M)$ 在某个开集 $W\subset M$ 上与 $\widetilde X,\widetilde Y\in\mathfrak{X}(M)$ 一致, 则 $\nabla_XY=\nabla_{\widetilde X}\widetilde Y$ 在 $W$ 上成立.
{\small (\textbf{提示:} 使用支集包含在 $W$ 中的隆起函数, 以及 Definition~\ref{def:affine-connection} 中列出的性质).}
\item 导出 $\nabla_XY$ 的局部坐标表达式 \eqref{eq:local-expression-of-affine-connection}.
\item 导出平行移动定律的局部坐标方程 \eqref{eq:parallel-transport-in-local-coordinates}.
\item 导出联络 $\nabla$ 的测地线的局部坐标方程 \eqref{eq:geodesic-equations-in-local-coordinates}.
\end{enumerate}
\item\label{ex:connections-with-straight-geodesics} 确定 $\mathbb{R}^n$ 上的所有仿射联络. 在这些联络中, 确定哪些联络的测地线是直线 $c(t)=at+b$ (其中 $a,b\in\mathbb{R}^n$).
\item\label{ex:covariant-derivative-of-tensors} 设 $\nabla$ 是 $M$ 上的仿射联络. 若 $\omega\in\Omega^1(M)$ 且 $X\in\mathfrak{X}(M)$, 则对所有 $Y\in\mathfrak{X}(M)$, 由
\[
  \nabla_X\omega(Y)=X\cdot(\omega(Y))-\omega(\nabla_XY)
\]
定义 $\omega$ 沿 $X$ 的\textbf{协变导数} $\nabla_X\omega\in\Omega^1(M)$.
\begin{enumerate}[label=(\alph*),leftmargin=*]
\item 证明此式确实定义了一个 1-形式, 即证明 $(\nabla_X\omega(Y))(p)$ 是 $Y_p$ 的线性函数.
\item 证明对所有 $X,Y\in\mathfrak{X}(M)$, $f,g\in C^\infty(M)$ 和 $\omega,\eta\in\Omega^1(M)$, 有
\begin{enumerate}[label=(\roman*),leftmargin=*]
\item $\nabla_{fX+gY}\omega=f\nabla_X\omega+g\nabla_Y\omega$;
\item $\nabla_X(\omega+\eta)=\nabla_X\omega+\nabla_X\eta$;
\item $\nabla_X(f\omega)=(X\cdot f)\omega+f\nabla_X\omega$.
\end{enumerate}
\item 设 $x:W\to\mathbb{R}^n$ 是开集 $W\subset M$ 上的局部坐标, 并取
\[
  \omega=\sum_{i=1}^n\omega_i\,\mathrm{d}x^i.
\]
证明
\[
  \nabla_X\omega=\sum_{i=1}^n
  \left(X\cdot\omega_i-\sum_{j,k=1}^n\Gamma_{ji}^kX^j\omega_k\right)\mathrm{d}x^i.
\]
\item 定义 $M$ 中任意张量场 $T$ 的协变导数 $\nabla_XT$, 并写出它在局部坐标中的表达式.
\end{enumerate}
\end{enumerate}
\end{RiemannExercise}
